\chapter{2007年重庆卷（文）}

\section{单选题}

\begin{problem}
在等比数列 \(\{a_{n}\}\) 中，\(a_{2} = 8\)，\(a_{1} = 64\)，则公比 \(q\) 为（\(\quad\)）

\choices
    {\(2\)}
    {\(3\)}
    {\(4\)}
    {\(8\)}
\end{problem}

\begin{answer}
\(\frac{1}{8}\)（按题面计算；原题四个选项均不含此值）
\end{answer}

\begin{solution}
等比数列满足 \(a_{2}=a_{1}q\)，所以 \(q=\dfrac{a_{2}}{a_{1}}=\dfrac{8}{64}=\dfrac{1}{8}\). 因此按题面计算公比为 \(\dfrac{1}{8}\)，但原题给出的四个选项均不包含该值.
\end{solution}

\begin{problem}
设全集 \( U = \{a, b, c, d\} \)，\( A = \{a, c\} \)，\( B = \{b\} \)，则 \( A \cap (\complement_U B) = \)（\(\quad\)）

\choices
    {\( \varnothing \)}
    {\( \{a\} \)}
    {\( \{c\} \)}
    {\( \{a, c\} \)}
\end{problem}

\begin{answer}
D
\end{answer}

\begin{solution}
全集中除去 \(B=\{b\}\) 后，\(\complement_U B=\{a,c,d\}\). 因此 \(A\cap(\complement_U B)=\{a,c\}\)，故选 D.
\end{solution}

\begin{problem}
垂直于同一平面的两条直线（\(\quad\)）

\choices
    {平行}
    {垂直}
    {相交}
    {异面}
\end{problem}

\begin{answer}
A
\end{answer}

\begin{solution}
设两条直线分别为 \(l_1\)，\(l_2\)，所给平面为 \(\alpha\). 两条直线都垂直于 \(\alpha\)，因而它们的方向都平行于平面 \(\alpha\) 的法线，故 \(l_1\parallel l_2\)，选 A.
\end{solution}

\begin{problem}
\((2x - 1)^{6}\) 展开式中 \(x^{2}\) 的系数为（\(\quad\)）

\choices
    {\(15\)}
    {\(60\)}
    {\(120\)}
    {\(240\)}
\end{problem}

\begin{answer}
B
\end{answer}

\begin{solution}
在二项式 \((2x-1)^6\) 中，含 \(x^2\) 的项须从六个因式中选出两个 \(2x\)，其系数为 \(\mathrm C_6^2\cdot 2^2\cdot(-1)^4=15\cdot4=60\)，故选 B.
\end{solution}

\begin{problem}
\(-1 < x < 1\) 是“\(x^2 < 1\)”的（\(\quad\)）

\choices
    {充分必要条件}
    {充分但不必要条件}
    {必要但不充分条件}
    {既不充分也不必要条件}
\end{problem}

\begin{answer}
A
\end{answer}

\begin{solution}
若 \(-1<x<1\)，则 \(|x|<1\)，从而 \(x^2<1\). 反之，若 \(x^2<1\)，则 \(|x|<1\)，即 \(-1<x<1\). 两个条件可以互相推出，所以前者是后者的充分必要条件，选 A.
\end{solution}

\begin{problem}
下列各式中，值为 \(\frac{\sqrt{3}}{2}\) 的是（\(\quad\)）

\choices
    {\(2\sin 15^{\circ}\cos 15^{\circ}\)}
    {\(\cos^2 15^\circ -\sin^2 15^\circ\)}
    {\(2\sin^{2}15^{\circ} - 1\)}
    {\(\sin^2 15^\circ +\cos^2 15^\circ\)}
\end{problem}

\begin{answer}
B
\end{answer}

\begin{solution}
利用三角恒等式分别计算四个选项，得
\(2\sin15^{\circ}\cos15^{\circ}=\sin30^{\circ}=\frac12\)，
\(\cos^2 15^{\circ}-\sin^2 15^{\circ}=\cos30^{\circ}=\frac{\sqrt3}{2}\)，
\(2\sin^2 15^{\circ}-1=-\cos30^{\circ}=-\frac{\sqrt3}{2}\)，
\(\sin^2 15^{\circ}+\cos^2 15^{\circ}=1\). 只有选项 B 的值符合要求，故选 B.
\end{solution}

\begin{problem}
从 \(5\) 张 \(100\) 元，\(3\) 张 \(200\) 元，\(2\) 张 \(300\) 元的奥运预赛门票中任取 \(3\) 张，则所取 \(3\) 张中至少有 \(2\) 张价格相同的概率为（\(\quad\)）

\choices
    {\( \frac{1}{4} \)}
    {\(\frac{79}{120}\)}
    {\( \frac{3}{4} \)}
    {\(\frac{23}{24}\)}
\end{problem}

\begin{answer}
C
\end{answer}

\begin{solution}
任取三张票的总数为 \(\mathrm C_{10}^{3}=120\). 事件“价格全不相同”要求分别取一张 \(100\) 元、\(200\) 元和 \(300\) 元的票，共有 \(5\cdot3\cdot2=30\) 种. 因此至少有两张价格相同的票共有 \(120-30=90\) 种，所求概率为 \(\dfrac{90}{120}=\dfrac34\)，故选 C.
\end{solution}

\begin{problem}
若直线 \( y = kx + 1 \) 与圆 \( x^{2} + y^{2} = 1 \) 相交于 \(P\)，\(Q\) 两点，且 \( \angle POQ = 120^{\circ} \) （其中 \(O\) 为原点），则 \( k \) 的值为（\(\quad\)）

\choices
    {\(-\sqrt{3}\) 或 \(\sqrt{3}\)}
    {\(\sqrt{3}\)}
    {\(-\sqrt{3}\) 或 \(\sqrt{2}\)}
    {\(\sqrt{2}\)}
\end{problem}

\begin{answer}
A
\end{answer}

\begin{solution}
圆的半径为 \(1\)，弦 \(PQ\) 所对的圆心角为 \(120^{\circ}\)，所以圆心到弦的距离为 \(\cos60^{\circ}=\frac12\). 直线 \(y=kx+1\) 写成 \(kx-y+1=0\)，原点到该直线的距离为 \(\dfrac{1}{\sqrt{k^2+1}}\). 由相交弦的条件得 \(\dfrac{1}{\sqrt{k^2+1}}=\frac12\)，即 \(k^2=3\)，所以 \(k=\pm\sqrt3\)，故选 A.
\end{solution}

\begin{problem}
已知向量 \(\overrightarrow{OA} = (4,6)\)，\(\overrightarrow{OB} = (3,5)\)，且 \(\overrightarrow{OC} \perp \overrightarrow{OA}\)，\(\overrightarrow{AC} \parallel \overrightarrow{OB}\)，则向量 \(\overrightarrow{OC} =\)（\(\quad\)）

\choices
    {\( \left(-\frac{3}{7},\frac{2}{7}\right) \)}
    {\( \left(-\frac{2}{7},\frac{4}{21}\right) \)}
    {\(\left(\frac{3}{7}, -\frac{2}{7}\right)\)}
    {\(\left(\frac{2}{7}, -\frac{4}{21}\right)\)}
\end{problem}

\begin{answer}
D
\end{answer}

\begin{solution}
设 \(\overrightarrow{OC}=(u,v)\). 由 \(\overrightarrow{AC}\parallel\overrightarrow{OB}\)，存在实数 \(t\) 使
\(\overrightarrow{OC}=\overrightarrow{OA}+t\overrightarrow{OB}=(4+3t,6+5t)\). 又因 \(\overrightarrow{OC}\perp\overrightarrow{OA}\)，有
\(4(4+3t)+6(6+5t)=0\)，解得 \(t=-\dfrac{26}{21}\). 于是
\(\overrightarrow{OC}=\left(\dfrac27,-\dfrac4{21}\right)\)，故选 D.
\end{solution}

\begin{problem}
设 \(P(3,1)\) 为二次函数 \(f(x) = ax^{2} - 2ax + b\)（\(x \geqslant 1\)）的图象与其反函数 \(y = f^{-1}(x)\) 的图象的一个交点. 则（\(\quad\)）

\choices
    {\(a = \frac{1}{2}, b = \frac{5}{2}\)}
    {\(a = \frac{1}{2}\)，\(b = -\frac{5}{2}\)}
    {\(a = -\frac{1}{2}\)，\(b = \frac{5}{2}\)}
    {\(a = -\frac{1}{2}\)，\(b = -\frac{5}{2}\)}
\end{problem}

\begin{answer}
C
\end{answer}

\begin{solution}
点 \(P(3,1)\) 在函数图象上，所以 \(f(3)=1\)，即 \(3a+b=1\). 点 \(P(3,1)\) 又在反函数图象上，按反函数关系有 \(f(1)=3\)，即 \(-a+b=3\). 两式相减得 \(4a=-2\)，从而 \(a=-\frac12\)，再得 \(b=\frac52\)，故选 C. 所得函数在 \(x\geqslant1\) 上单调递减，满足反函数存在所需的一一对应性.
\end{solution}

\begin{problem}
设 \(\sqrt{3}b\) 是 \(1 - a\) 和 \(1 + a\) 的等比中项，则 \(a + 3b\) 的最大值为（\(\quad\)）

\choices
    {\(1\)}
    {\(2\)}
    {\(3\)}
    {\(4\)}
\end{problem}

\begin{answer}
B
\end{answer}

\begin{solution}
由等比中项关系，\((\sqrt3b)^2=(1-a)(1+a)\)，所以 \(a^2+3b^2=1\). 由柯西不等式，
\(a+3b=a+\sqrt3(\sqrt3b)\leqslant\sqrt{a^2+3b^2}\sqrt{1+3}=2\). 当 \(a=\frac12\)，\(b=\frac12\) 时等号成立，故最大值为 \(2\)，选 B.
\end{solution}

\begin{problem}
已知以 \(F_{1}(-2,0)\)，\(F_{2}(2,0)\) 为焦点的椭圆与直线 \(x + \sqrt{3} y + 4 = 0\) 有且仅有一个交点，则椭圆的长轴长为（\(\quad\)）

\choices
    {\(3\sqrt{2}\)}
    {\(2\sqrt{6}\)}
    {\(2\sqrt{7}\)}
    {\(4\sqrt{2}\)}
\end{problem}

\begin{answer}
C
\end{answer}

\begin{solution}
设椭圆方程为 \(\dfrac{x^2}{a^2}+\dfrac{y^2}{b^2}=1\)，则 \(a>b>0\) 且 \(a^2-b^2=2^2=4\). 直线与椭圆恰有一个交点，故为椭圆的切线. 由柯西不等式，椭圆上 \(x+\sqrt3y\) 的最大值为 \(\sqrt{a^2+3b^2}\)，所以切线 \(x+\sqrt3y=-4\) 满足 \(a^2+3b^2=16\). 联立得 \(a^2+3(a^2-4)=16\)，即 \(a^2=7\). 椭圆长轴长为 \(2a=2\sqrt7\)，故选 C.
\end{solution}

\section{填空题}

\begin{problem}
在 \(\triangle ABC\) 中，\(AB = 1\)，\(BC = 2\)，\(B = 60^{\circ}\)，则 \(AC =\fillinblank{}\).
\end{problem}

\begin{answer}
\(\sqrt{3}\)
\end{answer}

\begin{solution}
在三角形 \(ABC\) 中，边 \(AC\) 对应角 \(B\). 由余弦定理，\(AC^2=AB^2+BC^2-2AB\cdot BC\cos B=1^2+2^2-2\cdot1\cdot2\cdot\cos60^{\circ}=3\). 因为 \(AC>0\)，所以 \(AC=\sqrt3\).
\end{solution}

\begin{problem}
已知 \(\begin{cases}
2x + 3y \leqslant 6,\\
x - y \geqslant 0,\\
y \geqslant 0,
\end{cases}\) 则 \(z = 3x - y\) 的最大值为\fillinblank{}.
\end{problem}

\begin{answer}
\(9\)
\end{answer}

\begin{solution}
由约束条件 \(y\geqslant0\) 且 \(x-y\geqslant0\)，可知 \(x\geqslant y\geqslant0\). 可行域的顶点为 \((0,0)\)、\((3,0)\) 以及直线 \(x=y\) 与 \(2x+3y=6\) 的交点 \(\left(\frac65,\frac65\right)\). 在三个顶点处，\(z=3x-y\) 的值分别为 \(0\)、\(9\)、\(\frac{12}{5}\)，所以最大值为 \(9\).
\end{solution}

\begin{problem}
要排出某班一天中语文，数学，政治，英语，体育，艺术\(6\)门课各一节的课程表，要求数学课排在前\(3\)节，英语课不排在第\(6\)节，则不同的排法种数为\fillinblank{}.（以数字作答）
\end{problem}

\begin{answer}
\(288\)
\end{answer}

\begin{solution}
先安排数学课. 数学课的位置有 \(3\) 种，其他五门课任意排列有 \(5!\) 种，故满足“数学课排在前 \(3\) 节”的排法有 \(3\cdot5!=360\) 种. 其中英语课排在第 \(6\) 节的排法有 \(3\cdot4!=72\) 种，因此所求排法数为 \(360-72=288\).
\end{solution}

\begin{problem}
函数 \( f(x)=\sqrt{x^{2}-2x}+2^{\sqrt{x^{2}-5x+4}} \) 的最小值为\fillinblank{}.
\end{problem}

\begin{answer}
\(1+2\sqrt{2}\)
\end{answer}

\begin{solution}
根式有意义须满足 \(x^2-2x\geqslant0\) 和 \(x^2-5x+4\geqslant0\)，所以定义域为 \(x\leqslant0\) 或 \(x\geqslant4\). 当 \(x\leqslant0\) 时，\(x^2-5x+4\geqslant4\)，从而 \(f(x)\geqslant2^2=4>1+2\sqrt2\). 当 \(x\geqslant4\) 时，\(\sqrt{x^2-2x}\geqslant\sqrt8=2\sqrt2\)，且 \(2^{\sqrt{x^2-5x+4}}\geqslant1\)，故 \(f(x)\geqslant1+2\sqrt2\). 在 \(x=4\) 时等号成立，所以函数的最小值为 \(1+2\sqrt2\).
\end{solution}

\section{解答题}

\begin{problem}
设甲，乙两人每次射击命中目标的概率分别为 \(\frac{3}{4}\) 和 \(\frac{4}{5}\)，且各次射击相互独立.

\begin{enumerate}
\item 若甲，乙各射击一次，求甲命中但乙未命中目标的概率；
\item 若甲，乙各射击两次，求两人命中目标的次数相等的概率.
\end{enumerate}
\end{problem}

\begin{solution}
\begin{enumerate}
\item 甲命中且乙未命中的概率为 \(\dfrac34\left(1-\dfrac45\right)=\dfrac34\cdot\dfrac15=\dfrac{3}{20}\).
\item 设甲、乙两人两次射击的命中次数分别为 \(X\)，\(Y\). 由二项分布和独立性，\(P(X=0)=\left(\frac14\right)^2=\frac1{16}\)，\(P(Y=0)=\left(\frac15\right)^2=\frac1{25}\)；\(P(X=1)=\mathrm C_2^1\frac34\frac14=\frac38\)，\(P(Y=1)=\mathrm C_2^1\frac45\frac15=\frac8{25}\)；\(P(X=2)=\left(\frac34\right)^2=\frac9{16}\)，\(P(Y=2)=\left(\frac45\right)^2=\frac{16}{25}\). 因此 \(P(X=Y)=\frac1{16}\cdot\frac1{25}+\frac38\cdot\frac8{25}+\frac9{16}\cdot\frac{16}{25}=\frac{1+48+144}{400}=\frac{193}{400}\).
\end{enumerate}
\end{solution}

\begin{problem}
已知函数 \( f(x) = \frac{1 + \sqrt{2}\cos\left(2x - \frac{\pi}{4}\right)}{\sin\left(x + \frac{\pi}{2}\right)} \).

\begin{enumerate}
\item 求 \( f(x) \) 的定义域.
\item 若角 \(\alpha\) 在第一象限且 \(\cos \alpha = \frac{3}{5}\)，求 \(f(\alpha)\).
\end{enumerate}
\end{problem}

\begin{solution}
\begin{enumerate}
\item 因为分母不能为零，且 \(\sin\left(x+\frac\pi2\right)=\cos x\)，所以 \(\cos x\ne0\)，从而定义域为 \(\left\{x\mid x\ne\frac\pi2+k\pi,\ k\in\Z\right\}\).
\item 由 \(\alpha\) 在第一象限且 \(\cos\alpha=\frac35\)，得 \(\sin\alpha=\frac45\). 又 \(\sqrt2\cos\left(2\alpha-\frac\pi4\right)=\cos2\alpha+\sin2\alpha\)，所以 \(f(\alpha)=\dfrac{1+\cos2\alpha+\sin2\alpha}{\cos\alpha}=\dfrac{1+\cos^2\alpha-\sin^2\alpha+2\sin\alpha\cos\alpha}{\cos\alpha}=2(\cos\alpha+\sin\alpha)=\frac{14}{5}\).
\end{enumerate}
\end{solution}

\begin{problem}
如图，在直三棱柱 \(ABC - A_{1}B_{1}C_{1}\) 中，\(\angle ABC = 90^{\circ}\)，\(AB = 1\)，\(BC = \frac{3}{2}\)，\(AA_{1} = 2\)；点 \(D\) 在棱 \(BB_{1}\) 上，\(BD = \frac{1}{3} BB_{1}\)；\(B_{1}E \perp A_{1}D\)，垂足为 \(E\). 求：

\begin{enumerate}
\item 求异面直线 \(A_{1}D\) 与 \(B_{1}C_{1}\) 的距离；
\item 四棱锥 \(C - ABDE\) 的体积.
\end{enumerate}

\bitmapfigure[max width=0.38\linewidth]{img/2007/chongqing_liberal/q19_fig1.png}
\end{problem}

\begin{solution}
\begin{enumerate}
\item 建立空间直角坐标系，令 \(B=(0,0,0)\)，\(A=(1,0,0)\)，\(C=\left(0,\frac32,0\right)\)，竖直方向为 \(z\) 轴. 于是 \(A_1=(1,0,2)\)，\(B_1=(0,0,2)\)，\(C_1=\left(0,\frac32,2\right)\)，\(D=\left(0,0,\frac23\right)\). 设 \(E=A_1+s(D-A_1)\)，则 \(E=\left(1-s,0,2-\frac43s\right)\). 由 \(B_1E\perp A_1D\)，有 \((E-B_1)\cdot(D-A_1)=0\)，即 \(-(1-s)+\frac{16}{9}s=0\)，解得 \(s=\frac9{25}\)，所以 \(E=\left(\frac{16}{25},0,\frac{38}{25}\right)\).

\(\overrightarrow{B_1E}=\left(\frac{16}{25},0,-\frac{12}{25}\right)\) 同时垂直于 \(A_1D\) 和 \(B_1C_1\)，故它是两异面直线的公垂线. 因此距离为 \(\lvert B_1E\rvert=\sqrt{\left(\frac{16}{25}\right)^2+\left(\frac{12}{25}\right)^2}=\frac45\).
\item 四点 \(A\)，\(B\)，\(D\)，\(E\) 都在平面 \(y=0\) 内，四边形 \(ABDE\) 在 \(xOz\) 平面上的坐标依次为 \((1,0)\)、\((0,0)\)、\((0,\frac23)\)、\((\frac{16}{25},\frac{38}{25})\). 由鞋带公式，\(S_{ABDE}=\frac12\left(\frac23\cdot\frac{16}{25}+\frac{38}{25}\right)=\frac{73}{75}\). 点 \(C\) 到平面 \(ABDE\) 的距离为 \(BC=\frac32\)，故 \(V_{C-ABDE}=\frac13S_{ABDE}\cdot\frac32=\frac13\cdot\frac{73}{75}\cdot\frac32=\frac{73}{150}\).
\end{enumerate}
\end{solution}

\begin{problem}
用长为 \(18 \mathrm{~m}\) 的钢条围成一个长方体形状的框架，要求长方体的长与宽之比为 \(2:1\)，问该长方体的长，宽，高各为多少时，其体积最大？ 最大体积是多少？
\end{problem}

\begin{solution}
设长方体的宽为 \(x\mathrm{~m}\)，则长为 \(2x\mathrm{~m}\). 框架共有十二条棱，所以 \(4(2x+x+h)=18\)，即 \(h=\frac92-3x\)，且 \(0<x<\frac32\). 体积为 \(V(x)=2x^2\left(\frac92-3x\right)=9x^2-6x^3\)，于是 \(V'(x)=18x(1-x)\). 因此 \(V\) 在 \(0<x<1\) 时递增，在 \(1<x<\frac32\) 时递减，最大值在 \(x=1\) 时取得. 此时长为 \(2\mathrm{~m}\)，宽为 \(1\mathrm{~m}\)，高为 \(\frac32\mathrm{~m}\)，最大体积为 \(2\cdot1\cdot1\cdot\frac32=3\mathrm{~m}^3\).
\end{solution}

\begin{problem}
如图，倾斜角为 \(\alpha\) 的直线经过抛物线 \(y^{2} = 8x\) 的焦点 \(F\)，且与抛物线交于 \(A\)，\(B\) 两点.

\begin{enumerate}
\item 求抛物线的焦点 \(F\) 的坐标及准线 \(l\) 方程；
\item 若 \(\alpha\) 为锐角，作线段 \(AB\) 的垂直平分线 \(m\) 交 \(x\) 轴于点 \(P\)，证明 \(|FP| - |FP| \cos 2\alpha\) 为定值，并求此定值.
\end{enumerate}

\FigureLayoutDeclare{img/2007/chongqing_liberal/q21_fig1.png}{9.0cm}{10bp 10bp 10bp 10bp}
\bitmapfigure[width=9.0cm]{img/2007/chongqing_liberal/q21_fig1.png}
\end{problem}

\begin{solution}
\begin{enumerate}
\item 抛物线 \(y^2=8x\) 与标准方程 \(y^2=2px\) 比较得 \(p=4\)，所以焦点为 \(F=(2,0)\)，准线为 \(l:x=-2\).
\item 令 \(t=\tan\alpha>0\)，过焦点的直线为 \(y=t(x-2)\). 它与抛物线的交点横坐标是方程 \(t^2(x-2)^2=8x\) 的两根，设为 \(x_A\)，\(x_B\). 由韦达定理，\(x_A+x_B=4+\frac8{t^2}\)，故弦 \(AB\) 的中点 \(M\) 满足 \(x_M=2+\frac4{t^2}\)，且 \(y_M=t(x_M-2)=\frac4t\).

\(AB\) 的斜率为 \(t\)，所以其垂直平分线过 \(M\)，斜率为 \(-\frac1t\). 令该直线与 \(x\) 轴交于 \(P\)，由 \(y_P=0\) 得 \(x_P=x_M+4=6+\frac4{t^2}\)，从而 \(\lvert FP\rvert=x_P-2=4+\frac4{t^2}=4\csc^2\alpha\). 因此 \(\lvert FP\rvert-\lvert FP\rvert\cos2\alpha=4\csc^2\alpha(1-\cos2\alpha)=4\csc^2\alpha\cdot2\sin^2\alpha=8\)，为定值 \(8\).
\end{enumerate}
\end{solution}

\begin{problem}
已知各项均为正数的数列 \(\{a_{n}\}\) 的前 \(n\) 项和 \(S_{n}\) 满足 \(S_{1} > 1\)，且 \(6S_{n} = (a_{n} + 1)(a_{n} + 2)\)，\(n \in \N_{+}\).

\begin{enumerate}
\item 求 \( \{a_{n}\} \) 的通项公式；
\item 设数列 \(\{b_n\}\) 满足 \(a_n(2^{b_n} - 1) = 1\)，并记 \(T_n\) 为 \(\{b_n\}\) 的前 \(n\) 项和，求证：\(3T_n + 1 > \log_2(a_n + 3)\)，\(n \in \N_{+}\).
\end{enumerate}
\end{problem}

\begin{solution}
\begin{enumerate}
\item 取 \(n=1\)，由 \(S_1=a_1>1\) 及题设得 \(6a_1=(a_1+1)(a_1+2)\)，即 \((a_1-1)(a_1-2)=0\)，所以 \(a_1=2\). 对任意 \(n\geqslant1\)，由 \(S_{n+1}=S_n+a_{n+1}\) 得 \(6a_{n+1}=(a_{n+1}+1)(a_{n+1}+2)-(a_n+1)(a_n+2)\). 展开并移项，得到 \(a_{n+1}^{2}-3a_{n+1}-a_n^{2}-3a_n=0\)，因而 \((a_{n+1}+a_n)(a_{n+1}-a_n-3)=0\). 因各项为正数，故 \(a_{n+1}=a_n+3\)，从而 \(a_n=2+3(n-1)=3n-1\).
\item 由 \(a_n(2^{b_n}-1)=1\)，得 \(2^{b_n}=\dfrac{a_n+1}{a_n}\). 令 \(R_n=\dfrac{2^{3T_n+1}}{a_n+3}\). 由第一问应有 \(a_n=3n-1\)，于是 \(R_1=\dfrac{2(3/2)^3}{5}=\dfrac{27}{20}>1\)，且
\(\frac{R_{n+1}}{R_n}=2^{3b_{n+1}}\frac{a_n+3}{a_{n+1}+3}=\left(\frac{3n+3}{3n+2}\right)^3\frac{3n+2}{3n+5}=\frac{(3n+3)^3}{(3n+2)^2(3n+5)}\). 令 \(u=3n+2>0\)，则分子与分母之差为 \((u+1)^3-u^2(u+3)=3u+1>0\)，所以 \(R_{n+1}>R_n\). 归纳得 \(R_n>1\)，即 \(2^{3T_n+1}>a_n+3\). 两边取以 \(2\) 为底的对数，得到 \(3T_n+1>\log_2(a_n+3)\).
\end{enumerate}
\end{solution}
